\chapter{Stringhe e array a più dimensioni}
\label{chap:stringhe}

Questo capitolo completa la lezione sugli array e i puntatori con due
applicazioni importantissime. Le \textbf{stringhe}, cioè i testi, in C sono
array di caratteri: vedremo come sono fatte, le funzioni della libreria
standard che le trattano e come si ordinano le parole di un testo con un
array di puntatori. Gli \textbf{array a più dimensioni}, come le matrici,
sono array di array: vedremo come stanno in memoria, come si passano alle
funzioni e come si allocano dinamicamente.

\section{Le stringhe}

\subsection{Un array di caratteri con un terminatore}

Le stringhe sono array a una dimensione di caratteri. In C una stringa è
\textbf{terminata dal carattere nullo} \code{'\textbackslash 0'}, un byte con
tutti i bit a zero, il cui valore decimale è 0: è così che le funzioni
riconoscono dove finisce il testo. Per questo la dimensione di una stringa
deve includere lo spazio per il terminatore: la stringa \code{"a"} ha
\textbf{due} elementi, \code{'a'} e \code{'\textbackslash 0'}. La
dichiarazione

\begin{lstlisting}[style=c]
char s[] = "abc";
\end{lstlisting}

è trattata dal compilatore come
\code{char s[] = \{'a', 'b', 'c', '\textbackslash 0'\};}: un array di 4
caratteri (figura~\ref{fig:stringa}). Il programmatore deve stare attento a
non superare la dimensione dell'array, perché il C non lo controlla.

\begin{figure}[H]
  \centering
  \begin{tikzpicture}[font=\FiraCodeNL\small]
    \celle[minimum width=1.0cm, minimum height=0.6cm]{s}{(0,0)}{a,b,c,\textbackslash 0}
    \foreach \k/\cod in {0/97, 1/98, 2/99, 3/0}
      \node[font=\FiraCodeNL\tiny, text=neutral!45!black, anchor=south] at ([yshift=1pt]s-\k.north) {\cod};
    \node[mnome, anchor=east] at ([xshift=-3pt]s-0.west) {s};
    \node[mnota, anchor=west, align=left] at ([xshift=0.4cm]s-3.east)
      {4 byte: i 3 caratteri e il terminatore\\(sopra, i codici ASCII; sotto, gli indici)};
  \end{tikzpicture}
  \caption{La stringa \code{"abc"} in memoria.}
  \label{fig:stringa}
\end{figure}

\subsection{Le stringhe costanti sono puntatori}

Una stringa costante, come \code{"abc"}, è trattata dal compilatore come un
\textbf{puntatore} al suo primo carattere. Così si può scrivere

\begin{lstlisting}[style=c]
char *p = "abc";

printf("%s %s\n", p, p + 1);
\end{lstlisting}

che stampa \code{abc bc}: la conversione \code{\%s} riceve un indirizzo e
stampa i caratteri da lì fino al \code{'\textbackslash 0'}, e \code{p + 1}
punta alla \code{b} (figura~\ref{fig:stringa-puntatore}). Per lo stesso motivo
\code{"abc"[1]} vale \code{'b'} e \code{*("abc" + 2)} vale \code{'c'}:
espressioni corrette, che però non vanno mai usate nel codice vero.

\begin{figure}[H]
  \centering
  \begin{tikzpicture}[font=\FiraCodeNL\small]
    \noindici
    \celle[minimum width=1.0cm, minimum height=0.6cm]{t}{(0,0)}{a,b,c,\textbackslash 0}
    \foreach \k/\testo in {0/p, 1/p + 1} {
      \node[font=\FiraCodeNL\scriptsize, text=mref, anchor=south] (l\k) at ([yshift=0.6cm]t-\k.north) {\testo};
      \draw[mrif] (l\k.south) -- (t-\k.north);
    }
    \draw[wtbrace] ([yshift=-3pt]t-0.south west) -- ([yshift=-3pt]t-2.south east);
    \node[mnota, anchor=north] at ([yshift=-0.32cm]t-1.south) {\code{\%s} con \code{p}: \code{abc}};
    \draw[wtbrace, draw=accent!70!black] ([yshift=-0.95cm]t-1.south west) -- ([yshift=-0.95cm]t-2.south east);
    \node[mnota, anchor=north, text=accent!70!black] at ([yshift=-1.25cm]$(t-1.south)!0.5!(t-2.south)$) {\code{\%s} con \code{p + 1}: \code{bc}};
  \end{tikzpicture}
  \caption{\code{\%s} stampa dall'indirizzo ricevuto fino al terminatore.}
  \label{fig:stringa-puntatore}
\end{figure}

\subsection{Array o puntatore?}

\code{char s[] = "abc";} e \code{char *p = "abc";} sembrano uguali, ma non lo
sono. Il primo crea un array di 4 caratteri, una copia della stringa che il
programma può modificare. Il secondo crea solo un puntatore alla stringa
costante, che sta in una zona di memoria \textbf{di sola lettura}: tentare di
modificarla è un comportamento indefinito, e su Linux il programma viene
fermato.

\begin{affianco}[0.5][c]
\begin{lstlisting}[style=c, name=modifica.c]
#include <stdio.h>

int main(void)
{
   char s[] = "abc";
   char *p = "abc";

   s[0] = 'x';
   printf("%s\n", s);
   p[0] = 'x';
   printf("%s\n", p);
   return 0;
}
\end{lstlisting}
\risultato
\begin{console}
$ ./modifica
xbc
Segmentation fault
\end{console}
\wtnota{\textit{Segmentation fault} (a volte seguito da \textit{core dumped})
è il messaggio della shell quando il sistema ferma un programma che ha
toccato memoria non sua.}
\end{affianco}

\subsection{La stringa vuota}

Il carattere nullo non va scritto in output: è un segnale per i programmi,
non un carattere da mostrare, e scriverlo in un file potrebbe confondere i
programmi che lo elaborano. Scrivere stringhe \textbf{vuote} (\code{""}, che
contiene solo il terminatore), invece, può essere utile. Per esempio, per
mettere il plurale solo quando serve:

\begin{affianco}[0.6][c]
\begin{lstlisting}[style=c]
for (nfrogs = 1; nfrogs <= 3; nfrogs++)
   printf("%d frog%s\n", nfrogs, (nfrogs == 1) ? "" : "s");
\end{lstlisting}
\risultato
\begin{console}
1 frog
2 frogs
3 frogs
\end{console}
\end{affianco}

\section{Contare le parole}

La funzione seguente conta le parole di una stringa, intese come sequenze di
caratteri separate da spazi. Il parametro \code{s} è un puntatore che
\textbf{scorre} la stringa: \code{++s} lo sposta sul carattere successivo, e
\code{*s} è il carattere su cui si trova. \code{isspace} (in \code{ctype.h})
è vera per gli spazi, le tabulazioni e gli a capo.

\begin{affianco}[0.58][c]
\begin{lstlisting}[style=c, name=word_cnt.c]
/* Count the number of words in a string. */

#include <stdio.h>
#include <ctype.h>

int word_cnt(const char *s)
{
   int    cnt = 0;

   while (*s != '\0') {
      while (isspace(*s))         /* skip white space */
         ++s;
      if (*s != '\0') {           /* found a word */
         ++cnt;
         while (!isspace(*s) && *s != '\0')
            ++s;                  /* skip the word */
      }
   }
   return cnt;
}

int main(void)
{
   printf("%d\n", word_cnt("  Il linguaggio   C  "));
   printf("%d\n", word_cnt(""));
   printf("%d\n", word_cnt("una"));
   return 0;
}
\end{lstlisting}
\risultato
\begin{console}
$ ./word_cnt
3
0
1
\end{console}
\end{affianco}

A ogni giro del ciclo esterno si saltano gli spazi; se non si è arrivati
alla fine, si è trovata una parola: si incrementa il contatore e si salta la
parola fino al prossimo spazio (figura~\ref{fig:word-cnt}).

\begin{figure}[H]
  \centering
  \begin{tikzpicture}[font=\FiraCodeNL\scriptsize,
      sp/.style={mcella, minimum width=0.42cm, minimum height=0.5cm, inner xsep=0pt, fill=neutral!10!white},
      ch/.style={mcella, minimum width=0.42cm, minimum height=0.5cm, inner xsep=0pt}]
    \noindici
    \celle[ch]{w}{(0,0)}{\char"2423,\char"2423,I,l,\char"2423,l,i,n,g,u,a,g,g,i,o,\char"2423,\char"2423,\char"2423,C,\char"2423,\char"2423,\textbackslash 0}
    \foreach \ia/\ib/\nparola in {2/3/1, 5/14/2, 18/18/3} {
      \draw[wtbrace, draw=mok] ([yshift=-3pt]w-\ia.south west) -- ([yshift=-3pt]w-\ib.south east);
      \node[font=\Lato\scriptsize\bfseries, text=mok, anchor=north] at ([yshift=-0.3cm]$(w-\ia.south)!0.5!(w-\ib.south)$) {parola \nparola};
    }
    \node[mnota, anchor=west, align=left] at ([xshift=0.3cm]w-21.east) {\code{\char"2423} sono spazi:\\vengono saltati};
  \end{tikzpicture}
  \caption{\code{word\_cnt} alterna due fasi: salta gli spazi, poi salta una
           parola contandola.}
  \label{fig:word-cnt}
\end{figure}

\section{La parola chiave \code{const}}

\subsection{Variabili che non cambiano}

La parola chiave \code{const}, introdotta dall'ANSI C, indica che una
variabile può essere inizializzata ma non modificata
(capitolo~\ref{chap:primi-programmi}). Con \code{const int k = 3;}
l'istruzione \code{k = 1;} è un errore: \textit{assignment of read-only
variable ‘k’}.

Con i puntatori, \code{const} può riguardare il dato puntato o il puntatore
stesso:

\begin{itemize}
  \item \code{const int *p;} dichiara un puntatore a un \code{const int}:
        il \textbf{dato puntato} non si può modificare attraverso \code{p}
        (\code{*p = 1;} è un errore), ma il puntatore sì (\code{p = \&j;} va
        bene);
  \item \code{int *const p = \&i;} è il contrario: un puntatore
        \textbf{costante}, che punterà sempre a \code{i}, attraverso cui
        però \code{i} si può modificare.
\end{itemize}

\subsection{Una promessa verificata dal compilatore}

Il parametro di \code{word\_cnt} è \code{const char *s}: la funzione
\textbf{promette} di non modificare la stringa che riceve (sposta \code{s},
ma non tocca i caratteri). È un'informazione preziosa per chi la usa, e il
compilatore la fa rispettare. Se nella funzione si prova a scrivere nella
stringa:

\begin{lstlisting}[style=c, name=word_cnt2.c]
/* Count the number of words in a string. */

#include <ctype.h>

int word_cnt(const char *s)
{
   int    cnt = 0;

   while (*s != '\0') {
      while (isspace(*s))         /* skip white space */
         ++s;
      if (*s != '\0') {           /* found a word */
         ++cnt;
         while (!isspace(*s) && *s != '\0') {
            *s = 'a';             /* the compiler will complain */
            ++s;                  /* skip the word */
         }
      }
   }
   return cnt;
}
\end{lstlisting}

\begin{console}
$ gcc -Wall -c word_cnt2.c
word_cnt2.c: In function (*@‘@*)word_cnt(*@’@*):
word_cnt2.c:15:16: error: assignment of read-only location (*@‘@*)*s(*@’@*)
   15 |             *s = 'a';             /* the compiler will complain */
      |                ^
\end{console}

\section{Funzioni per caratteri e stringhe}

\subsection{I caratteri: \code{ctype.h}}

La libreria standard ha molte funzioni per classificare e trasformare i
caratteri, con i prototipi in \code{ctype.h}. Ricevono un \code{int} (come
\code{getchar}, capitolo~\ref{chap:tipi}) e le funzioni di verifica
restituiscono un valore diverso da zero (vero) o zero (falso).

\begin{center}
\renewcommand{\arraystretch}{1.12}
\begin{tabular}{@{}>{\FiraCodeNL}l l@{}}
  \toprule
  \normalfont\textbf{funzione} & \textbf{che cosa fa} \\
  \midrule
  int isblank(int c) & vero se \code{c} è uno spazio o una tabulazione (non c'era nei vecchi compilatori Microsoft) \\
  int isdigit(int c) & vero se \code{c} è una cifra \\
  int isalpha(int c) & vero se \code{c} è una lettera \\
  int isalnum(int c) & vero se \code{c} è una lettera o una cifra \\
  int islower(int c) & vero se \code{c} è una lettera minuscola \\
  int isupper(int c) & vero se \code{c} è una lettera maiuscola \\
  int tolower(int c) & se \code{c} è maiuscola restituisce la minuscola, altrimenti \code{c} \\
  int toupper(int c) & se \code{c} è minuscola restituisce la maiuscola, altrimenti \code{c} \\
  \bottomrule
\end{tabular}
\end{center}

\begin{lstlisting}[style=c, name=ctipi.c]
#include <stdio.h>
#include <ctype.h>

int main(void)
{
   const char *s = "aZ 7!";
   int i;

   printf("car  isalpha isdigit isalnum islower isupper isblank toupper\n");
   for (i = 0; s[i] != '\0'; i++)
      printf(" '%c' %5d %7d %7d %7d %7d %7d %7c\n", s[i],
             isalpha(s[i]) != 0, isdigit(s[i]) != 0, isalnum(s[i]) != 0,
             islower(s[i]) != 0, isupper(s[i]) != 0, isblank(s[i]) != 0,
             toupper(s[i]));
   return 0;
}
\end{lstlisting}

\begin{console}
$ ./ctipi
car  isalpha isdigit isalnum islower isupper isblank toupper
 'a'     1       0       1       1       0       0       A
 'Z'     1       0       1       0       1       0       Z
 ' '     0       0       0       0       0       1
 '7'     0       1       1       0       0       0       7
 '!'     0       0       0       0       0       0       !
\end{console}

Nel programma, \code{!= 0} trasforma qualunque “vero” in 1, per avere una
tabella leggibile.

\subsection{Le stringhe: \code{string.h}}

Le funzioni per le stringhe hanno i prototipi in \code{string.h}. Il
parametro \code{const char *} indica una stringa che la funzione non
modifica; il \code{char *} senza \code{const} è l'array in cui la funzione
scrive, che \textbf{deve avere spazio sufficiente}: nessuna di queste
funzioni lo controlla.

\begin{center}
\renewcommand{\arraystretch}{1.12}
\begin{tabular}{@{}>{\FiraCodeNL\small}l>{\raggedright\arraybackslash}p{8.4cm}@{}}
  \toprule
  \normalfont\textbf{prototipo} & \textbf{che cosa fa} \\
  \midrule
  size\_t strlen(const char *s) & la lunghezza di \code{s}, senza contare il terminatore \\
  char *strcpy(char *s1, const char *s2) & copia \code{s2} nell'array \code{s1}; restituisce \code{s1} \\
  char *strncpy(char *s1, const char *s2, size\_t n) & copia al massimo \code{n} caratteri di \code{s2} in \code{s1} \\
  char *strcat(char *s1, const char *s2) & aggiunge \code{s2} in coda a \code{s1}; restituisce \code{s1} \\
  char *strncat(char *s1, const char *s2, size\_t n) & aggiunge in coda al massimo \code{n} caratteri di \code{s2} \\
  int strcmp(const char *s1, const char *s2) & confronta: restituisce 0 se sono uguali, un negativo se \code{s1} viene prima, un positivo se viene dopo \\
  int strncmp(const char *s1, const char *s2, size\_t n) & come \code{strcmp}, ma solo sui primi \code{n} caratteri \\
  char *strchr(const char *s, int c) & un puntatore alla prima \code{c} in \code{s}, o \code{NULL} \\
  char *strrchr(const char *s, int c) & un puntatore all'ultima \code{c} in \code{s}, o \code{NULL} \\
  char *strstr(const char *s1, const char *s2) & un puntatore alla prima occorrenza di \code{s2} in \code{s1}, o \code{NULL} \\
  char *strpbrk(const char *s1, const char *s2) & un puntatore al primo carattere di \code{s1} che compare in \code{s2}, o \code{NULL} \\
  size\_t strspn(const char *s1, const char *s2) & la lunghezza del tratto iniziale di \code{s1} fatto \textbf{solo} di caratteri di \code{s2} \\
  size\_t strcspn(const char *s1, const char *s2) & la lunghezza del tratto iniziale di \code{s1} fatto di caratteri \textbf{non} in \code{s2} \\
  \bottomrule
\end{tabular}
\end{center}

\code{strcmp} confronta i codici dei caratteri uno alla volta, finché trova
una differenza: per questo l'ordine è quello alfabetico per le minuscole (o
per le maiuscole), ma tutte le maiuscole vengono prima delle minuscole
(capitolo~\ref{chap:tipi}). Ecco gli esempi delle slide per le ultime due
funzioni:

\begin{affianco}[0.55][c]
\begin{lstlisting}[style=c, name=spn.c]
printf("%zu\n", strcspn("abc", "abc"));       /* 0 */
printf("%zu\n", strcspn("abc", "abcde"));     /* 0 */
printf("%zu\n", strcspn("djpc", "lqrabc"));   /* 3 */
printf("%zu\n", strspn("abc", "abc"));        /* 3 */
printf("%zu\n", strspn("abc", "abcde"));      /* 3 */
printf("%zu\n", strspn("djpc", "lqrabc"));    /* 0 */
\end{lstlisting}
\risultato
\begin{itemize}
  \item \code{strcspn("djpc", "lqrabc")} vale 3: \code{d}, \code{j} e
        \code{p} non sono in \code{"lqrabc"}, la \code{c} sì.
  \item \code{strspn("abc", "abcde")} vale 3: tutti i caratteri di
        \code{"abc"} sono in \code{"abcde"}.
  \item I valori sono di tipo \code{size\_t}, che si stampa con \code{\%zu}
        (nelle slide c'è \code{\%d}, e \code{gcc} avvisa).
\end{itemize}
\end{affianco}

\subsection{Come sono fatte}

Non c'è niente di speciale in queste funzioni: si possono scrivere in C, in
poche righe. Ecco le versioni delle slide (con il prefisso \code{my\_}: i
nomi della libreria standard vanno trattati come riservati,
capitolo~\ref{chap:sintassi}):

\begin{affianco}[0.5][c]
\begin{lstlisting}[style=c, name=mystr.c]
size_t my_strlen(const char *s)
{
   size_t n;

   for (n = 0; *s != '\0'; ++s)
      ++n;
   return n;
}

char *my_strcpy(char *s1, const char *s2)
{
   char *p = s1;

   while ((*p++ = *s2++))
      ;
   return s1;
}

char *my_strcat(char *s1, const char *s2)
{
   char *p = s1;

   while (*p)
      ++p;
   while ((*p++ = *s2++))
      ;
   return s1;
}
\end{lstlisting}
\risultato
\begin{lstlisting}[style=c]
int main(void)
{
   char *s1;

   s1 = malloc(10 * sizeof(char));
   my_strcpy(s1, "abc");
   printf("%s %zu\n", s1, my_strlen(s1));
   my_strcat(s1, "efg");
   printf("%s %zu\n", s1, my_strlen(s1));
   free(s1);
   return 0;
}
\end{lstlisting}
\begin{console}
$ ./mystr
abc 3
abcefg 6
\end{console}
\end{affianco}

Il cuore di \code{my\_strcpy} è il ciclo \code{while ((*p++ = *s2++))}, che
fa tutto nella condizione (figura~\ref{fig:strcpy}): copia il carattere
puntato da \code{s2} nella cella puntata da \code{p}, sposta entrambi i
puntatori in avanti (\code{++} postfisso: si usa il valore e poi si
incrementa), e il valore dell'assegnazione, cioè il carattere copiato,
fa da condizione. Quando si copia il \code{'\textbackslash 0'}, che vale 0,
il ciclo si ferma: il terminatore è stato copiato anche lui. Le doppie
parentesi dicono al compilatore (e al lettore) che l'assegnazione nella
condizione è voluta: le slide hanno le parentesi singole, e \code{gcc -Wall}
suggerirebbe di aggiungerle. \code{my\_strcat} prima porta \code{p} sul
terminatore di \code{s1}, poi copia da lì.

\begin{figure}[H]
  \centering
  \begin{tikzpicture}[font=\FiraCodeNL\small]
    \noindici
    \celle[minimum width=0.9cm, minimum height=0.55cm]{src}{(0,0)}{a,b,c,\textbackslash 0}
    \celle[minimum width=0.9cm, minimum height=0.55cm, fill=accent!8!white]{dst}{(0,-1.6)}{a,b,c,\textbackslash 0,\mvuota,\mvuota}
    \node[mnome, anchor=east] at ([xshift=-3pt]src-0.west) {s2};
    \node[mnome, anchor=east] at ([xshift=-3pt]dst-0.west) {s1};
    \foreach \k in {0,1,2,3}
      \draw[wtarrowacc] (src-\k.south) -- node[right, font=\Lato\tiny, text=accent!70!black] {\the\numexpr\k+1\relax} (dst-\k.north);
    \node[mnota, anchor=west, align=left] at ([xshift=0.5cm]src-3.east)
      {un giro per carattere: \code{*p = *s2}, poi\\\code{p} e \code{s2} avanzano; al giro 4 si copia\\il \code{'\textbackslash 0'}, che vale 0, e il ciclo finisce};
  \end{tikzpicture}
  \caption{\code{my\_strcpy(s1, s2)}: quattro giri, l'ultimo copia il
           terminatore.}
  \label{fig:strcpy}
\end{figure}

\subsection{Lo spazio per la copia}

Nelle slide, il programma di prova di \code{strcpy} riserva per \code{s1}
\code{malloc(3*sizeof(char))}, cioè 3 byte, e poi ci copia \code{"abc"}, che
ne occupa 4 con il terminatore. È proprio l'errore da cui mettono in guardia
le funzioni di \code{string.h}: si scrive un byte oltre la memoria riservata.
Il programma sembra funzionare, ma \code{gcc} se ne accorge già in
compilazione, e \code{-fsanitize=address} lo conferma in esecuzione:

\begin{console}
$ gcc -Wall -o strcpy3 strcpy3.c
strcpy3.c: In function (*@‘@*)main(*@’@*):
strcpy3.c:11:8: warning: (*@‘@*)strcpy(*@’@*) writing 4 bytes into a region of size 3 overflows the destination [-Wstringop-overflow=]
   11 |   s1 = strcpy(s1, s2);
      |        ^~~~~~~~~~~~~~
$ ./strcpy3
abc abc
$ gcc -g -fsanitize=address -o strcpy3 strcpy3.c
$ ./strcpy3
=================================================================
==3088==ERROR: AddressSanitizer: heap-buffer-overflow on address 0x502000000013 at pc 0x7f5d792a7923 bp 0x7ffec73255a0 sp 0x7ffec7324d48
WRITE of size 4 at 0x502000000013 thread T0
(*@\textit{\dots}@*)
\end{console}

La regola: per copiare una stringa \code{s2} servono
\code{strlen(s2) + 1} byte, cioè \code{malloc(strlen(s2) + 1)}.

\section{Gli array di puntatori}

\subsection{Un array di stringhe}

La dichiarazione \code{int *a[20];} definisce \code{a} come un array di 20
elementi, ognuno dei quali è un \textbf{puntatore a intero}: può essere
pensato come un array di vettori, ciascuno con la sua lunghezza. Il caso più
comune è l'array di stringhe, \code{char *w[N];}: ogni \code{w[i]} punta a
una stringa, e le stringhe possono essere lunghe quanto serve
(figura~\ref{fig:array-puntatori}).

\begin{figure}[H]
  \centering
  \begin{tikzpicture}[font=\FiraCodeNL\scriptsize,
      cc/.style={mcella, minimum width=0.45cm, minimum height=0.5cm, inner xsep=0pt}]
    \noindici
    \memoria[minimum width=1.0cm]{w}{(0,0)}{w[0]/{}/, w[1]/{}/, w[2]/{}/, w[3]/{}/}
    \celle[cc]{s0}{(3.2,0)}{l,e,\textbackslash 0}
    \celle[cc]{s1}{(3.2,-0.52)}{p,a,r,o,l,e,\textbackslash 0}
    \celle[cc]{s2}{(3.2,-1.04)}{d,i,\textbackslash 0}
    \celle[cc]{s3}{(3.2,-1.56)}{u,n,\textbackslash 0}
    \foreach \k [evaluate=\k as \r using int(\k+1)] in {0,1,2,3}
      \riferimento{w-\r}{-- (s\k-0.west)}
    \node[mnota, anchor=west, align=left] at (7.2,-0.8) {ogni stringa occupa\\solo lo spazio che le serve};
  \end{tikzpicture}
  \caption{Un array di puntatori a stringhe di lunghezza diversa.}
  \label{fig:array-puntatori}
\end{figure}

\subsection{Ordinare le parole di un testo}

L'esempio delle slide è un programma, diviso in più file
(capitolo~\ref{chap:make}), che legge le parole dallo standard input e le
stampa in ordine lessicografico (alfabetico, secondo i codici ASCII).

\begin{affianco}[0.5][t]
\codicepiccolo
\begin{lstlisting}[style=h, name=sort.h]
#define   MAXWORD   50        /* max word size */
#define   N         300       /* array size of w[] */

void   error_exit_calloc_failed(void);
void   error_exit_too_many_words(void);
void   error_exit_word_too_long(void);
void   sort_words(char *w[], int n);
void   swap(char **p, char **q);
void   wrt_words(char *w[], int n);
\end{lstlisting}
\begin{lstlisting}[style=c, name=sort_words.c]
#include <string.h>
#include "sort.h"

void sort_words(char *w[], int n) /* n elements are to be sorted */
{
   int   i, j;

   for (i = 0; i < n; ++i)
      for (j = i + 1; j < n; ++j)
         if (strcmp(w[i], w[j]) > 0)
            swap(&w[i], &w[j]);
}
\end{lstlisting}
\begin{lstlisting}[style=c, name=swap.c]
#include "sort.h"

void swap(char **p, char **q)
{
   char   *tmp;

   tmp = *p;
   *p = *q;
   *q = tmp;
}
\end{lstlisting}
\risultato
\begin{lstlisting}[style=c, name=main.c]
/* Sort words lexicographically. */
#include <stdio.h>
#include <stdlib.h>
#include <string.h>
#include "sort.h"

int main(void)
{
   char   word[MAXWORD];      /* work space */
   char   *w[N];              /* an array of pointers */
   int    n;                  /* number of words to be sorted */
   int    i;

   for (i = 0; scanf("%s", word) == 1; ++i) {
      if (i >= N)
         error_exit_too_many_words();
      if (strlen(word) >= MAXWORD)
         error_exit_word_too_long();
      w[i] = calloc(strlen(word) + 1, sizeof(char));
      if (w[i] == NULL)
         error_exit_calloc_failed();
      strcpy(w[i], word);
   }
   n = i;
   sort_words(w, n);          /* sort the words */
   wrt_words(w, n);           /* write sorted list of words */
   return 0;
}
\end{lstlisting}
\begin{lstlisting}[style=c, name=wrt.c]
#include <stdio.h>
#include "sort.h"

void wrt_words(char *w[], int n)
{
   int   i;

   for (i = 0; i < n; ++i)
      printf("%s\n", w[i]);
}
\end{lstlisting}
\end{affianco}

\subsection{Come funziona}

\begin{itemize}
  \item \code{scanf("\%s", word)} legge una parola (una sequenza di
        caratteri senza spazi) nell'array \code{word}, aggiungendo il
        terminatore. Il ciclo continua finché \code{scanf} riesce a leggere.
  \item Per ogni parola si controlla che non ce ne siano troppe e che non
        sia troppo lunga; poi \code{calloc} riserva esattamente
        \code{strlen(word) + 1} caratteri, e \code{strcpy} ci copia la
        parola. \code{w[i]} punta alla copia: \code{word} verrà riusato per
        la parola successiva.
  \item \code{sort\_words} è un ordinamento per selezione con \code{strcmp}:
        se \code{w[i]} viene dopo \code{w[j]}, si scambiano.
  \item \code{swap} riceve gli \textbf{indirizzi} di due elementi di
        \code{w}, cioè dei \code{char **}, e scambia i due
        \textbf{puntatori}: le stringhe non si spostano, cambia solo
        l'ordine delle frecce (figura~\ref{fig:scambio-puntatori}).
  \item Il file \code{error.c} contiene le tre funzioni di errore, che
        stampano un messaggio e terminano il programma con \code{exit(1)}.
\end{itemize}

\begin{figure}[H]
  \centering
  \begin{tikzpicture}[font=\FiraCodeNL\scriptsize,
      cc/.style={mcella, minimum width=0.45cm, minimum height=0.5cm, inner xsep=0pt}]
    \noindici
    \node[font=\Lato\scriptsize\bfseries] at (2.2,0.75) {prima dello scambio};
    \memoria[minimum width=1.0cm]{a}{(0,0)}{w[0]/{}/, w[1]/{}/}
    \celle[cc]{x0}{(2.6,0)}{p,e,r,a,\textbackslash 0}
    \celle[cc]{x1}{(2.6,-0.52)}{m,e,l,a,\textbackslash 0}
    \riferimento{a-1}{-- (x0-0.west)}
    \riferimento{a-2}{-- (x1-0.west)}
    \begin{scope}[xshift=7.6cm]
      \node[font=\Lato\scriptsize\bfseries] at (2.2,0.75) {dopo \code{swap(\&w[0], \&w[1])}};
      \memoria[minimum width=1.0cm]{b}{(0,0)}{w[0]/{}/, w[1]/{}/}
      \celle[cc]{y0}{(2.6,0)}{p,e,r,a,\textbackslash 0}
      \celle[cc]{y1}{(2.6,-0.52)}{m,e,l,a,\textbackslash 0}
      \riferimento[draw=accent!75!black]{b-1}{to[out=0, in=180] (y1-0.west)}
      \riferimento[draw=accent!75!black]{b-2}{to[out=0, in=180] (y0-0.west)}
    \end{scope}
  \end{tikzpicture}
  \caption{\code{swap} scambia i puntatori, non i caratteri: ora
           \code{w[0]} punta a \code{"mela"}.}
  \label{fig:scambio-puntatori}
\end{figure}

\subsection{Il programma in esecuzione}

Si compilano insieme i cinque file sorgente e si prova con una frase:

\begin{console}
$ gcc -Wall -o sort_words main.c error.c wrt.c swap.c sort_words.c
$ echo "le parole di un Testo da Ordinare" | ./sort_words
Ordinare
Testo
da
di
le
parole
un
\end{console}

\code{echo} scrive la frase sullo standard output, e il simbolo \code{|}
(\textit{pipe}) la collega allo standard input di \code{sort\_words}: è un
altro modo di fornire l'input, oltre alla tastiera e alla redirezione
\code{<}. Le parole con la maiuscola vengono prima di tutte le altre, perché
in ASCII le maiuscole precedono le minuscole.

Nella cartella degli esempi c'è anche un file \code{input.txt} con una frase
di 50 parole (una citazione sugli array a lunghezza variabile, che vedremo
alla fine del capitolo); \code{./sort\_words < input.txt} le stampa
ordinate, da \code{According} e \code{Arrays} fino a \code{world}.

\section{Gli array a più dimensioni}

\subsection{Array di array}

In C si possono definire array di array:

\begin{lstlisting}[style=c]
int a[100];           /* array a una dimensione */
int b[2][7];          /* array a due dimensioni */
int c[5][3][2];       /* array a tre dimensioni */
\end{lstlisting}

Tutti gli elementi di un array sono memorizzati in memoria
\textbf{contigua}, a partire dal primo: \code{a[0]}, \code{b[0][0]},
\code{c[0][0][0]}. Un array a due dimensioni si può pensare come una
\textbf{matrice}: con \code{int a[3][5];}, 3 è il numero di righe e 5 il
numero di colonne. In memoria, però, la matrice è “srotolata”: prima tutta
la riga 0, poi la riga 1, e così via (figura~\ref{fig:matrice}).

\begin{affianco}[0.55][c]
\begin{lstlisting}[style=c, name=matrice.c]
#include <stdio.h>

int main(void)
{
   int a[3][5], i, j;

   for (i = 0; i < 3; i++)
      for (j = 0; j < 5; j++)
         a[i][j] = 10 * i + j;

   for (i = 0; i < 3; i++) {
      for (j = 0; j < 5; j++)
         printf("%4d", a[i][j]);
      printf("\n");
   }
   printf("&a[0][0] = %p\n", (void *) &a[0][0]);
   printf("&a[0][4] = %p\n", (void *) &a[0][4]);
   printf("&a[1][0] = %p\n", (void *) &a[1][0]);
   i = 2; j = 3;
   printf("%d %d %d %d %d\n", a[i][j], *(a[i] + j), (*(a + i))[j],
          *((*(a + i)) + j), *(&a[0][0] + 5 * i + j));
   return 0;
}
\end{lstlisting}
\risultato
\begin{console}
$ ./matrice
   0   1   2   3   4
  10  11  12  13  14
  20  21  22  23  24
&a[0][0] = 0x7ffff4d0aaa0
&a[0][4] = 0x7ffff4d0aab0
&a[1][0] = 0x7ffff4d0aab4
23 23 23 23 23
\end{console}
\wtnota{\code{a[1][0]} sta 4 byte dopo \code{a[0][4]}: subito dopo
l'ultimo elemento della riga 0.}
\end{affianco}

\begin{figure}[H]
  \centering
  \begin{tikzpicture}[font=\FiraCodeNL\scriptsize,
      r0/.style={mcella, minimum width=0.62cm},
      r1/.style={mcella, minimum width=0.62cm, fill=accent!10!white},
      r2/.style={mcella, minimum width=0.62cm, fill=green-gradient-1!40!white}]
    \noindici
    \node[font=\Lato\scriptsize\bfseries] at (1.6,0.75) {come la pensiamo};
    \celle[r0]{m0}{(0,0)}{0,1,2,3,4}
    \celle[r1]{m1}{(0,-0.5)}{10,11,12,13,14}
    \celle[r2]{m2}{(0,-1.0)}{20,21,22,23,24}
    \foreach \k in {0,1,2}
      \node[mindir, anchor=east] at ([xshift=-2pt]m\k-0.west) {riga \k};
    \node[font=\Lato\scriptsize\bfseries] at (10.0,0.75) {come sta in memoria};
    \celle[r0]{l0}{(5.0,-0.5)}{0,1,2,3,4}
    \celle[r1]{l1}{(l0-4.east)}{10,11,12,13,14}
    \celle[r2]{l2}{(l1-4.east)}{20,21,22,23,24}
    \node[mnota, anchor=north, align=center] at ([yshift=-0.2cm]l1-2.south) {\code{a[i][j]} sta a $5 \cdot i + j$ posti da \code{a[0][0]}};
  \end{tikzpicture}
  \caption{La matrice \code{int a[3][5]}: righe una dopo l'altra in memoria.}
  \label{fig:matrice}
\end{figure}

\subsection{La storage mapping function}

La funzione che, per ogni array, associa agli indici la posizione
dell'elemento si chiama \term{storage mapping function}. Per un array a una
dimensione è \code{a[i]} $\rightarrow$ \code{*(\&a[0] + i)}. Per un array a
due dimensioni, con $N$ colonne,

\[
  \code{a[i][j]} \;\rightarrow\; \code{*(\&a[0][0] + } N \cdot i + j\code{)}
\]

perché per arrivare alla riga $i$ bisogna saltare $i$ righe intere da $N$
elementi. Il numero di righe (la prima dimensione) \textbf{non viene usato}.
Per lo stesso motivo \code{a[i][j]} si può scrivere anche
\code{*(a[i] + j)}, \code{(*(a + i))[j]} o \code{*((*(a + i)) + j)}: l'ultima
riga di \code{matrice.c} li stampa tutti, e valgono tutti 23.

Lo stesso vale con più dimensioni. Con \code{int A[7][9][2];} il compilatore
riserva spazio contiguo per $7 \cdot 9 \cdot 2 = 126$ interi, l'indirizzo
base è \code{\&A[0][0][0]}, e

\[
  \code{A[i][j][k]} \;\rightarrow\; \code{*(\&A[0][0][0] + } 9 \cdot 2 \cdot i + 2 \cdot j + k\code{)}
\]

Anche qui la prima dimensione, 7, non serve.

\subsection{L'inizializzazione}

Un array a più dimensioni si può inizializzare in molti modi. Queste
inizializzazioni sono equivalenti:

\begin{lstlisting}[style=c]
int A[2][3] = {1, 2, 3, 4, 5, 6};
int A[2][3] = {{1, 2, 3}, {4, 5, 6}};
int A[][3] = {{1, 2, 3}, {4, 5, 6}};
int A[][3] = {1, 2, 3, 4, 5, 6};
\end{lstlisting}

\begin{itemize}
  \item Senza graffe interne, gli elementi vengono riempiti per riga:
        \code{A[0][0]}, \code{A[0][1]}, \code{A[0][2]}, \code{A[1][0]},
        \dots{} (\code{gcc -Wall} suggerisce comunque le graffe:
        \textit{missing braces around initializer}).
  \item Se i valori sono meno degli elementi, i rimanenti valgono 0:
        \code{\{\{1\}, \{4, 5\}\}} dà le righe \code{1 0 0} e \code{4 5 0}.
  \item Se la prima dimensione è vuota, il compilatore la ricava dal numero
        di graffe interne (o dei valori); tutte le altre dimensioni vanno
        indicate.
  \item \code{int A[2][2][3] = \{0\};} azzera tutto l'array; qui vanno
        dichiarate tutte le dimensioni. È il modo più semplice di inizializzare
        un array automatico, che altrimenti conterrebbe valori a caso.
\end{itemize}

Anche queste due inizializzazioni di un array a tre dimensioni sono
equivalenti (verificato stampando gli elementi: \code{1 1 0 2 0 0 3 0 0 4 4
0}):

\begin{lstlisting}[style=c]
int A[2][2][3] = {
    {{1, 1, 0}, {2, 0, 0}},
    {{3, 0, 0}, {4, 4, 0}}
};
int A[][2][3] = {{{1, 1}, {2}}, {{3}, {4, 4}}};
\end{lstlisting}

\subsection{Le matrici come parametri}

Una matrice si può passare a una funzione, ma per generare il codice il
compilatore ha bisogno della storage mapping function: nella dichiarazione
del parametro vanno quindi indicate \textbf{tutte le dimensioni tranne la
prima}. Queste tre intestazioni sono equivalenti:

\begin{lstlisting}[style=c]
int sum(int a[3][5])
int sum(int a[][5])
int sum(int (*a)[5])        /* a punta ad array di 5 int */
\end{lstlisting}

ma, se il corpo usa \code{i < 3} e \code{j < 5}, vanno bene solo per le
matrici $3 \times 5$. Passando il numero di righe come parametro, la funzione
va bene per tutte le matrici $m \times 5$:

\begin{affianco}[0.5][c]
\begin{lstlisting}[style=c, name=sum2d.c]
int sum(int a[][5], int m)
{
   int   i, j, sum = 0;

   for (i = 0; i < m; ++i)
      for (j = 0; j < 5; ++j)
         sum += a[i][j];
   return sum;
}
\end{lstlisting}
\risultato
Il numero di colonne, però, resta fisso. Come scrivere una funzione per
matrici di qualunque dimensione? La prima idea, \code{int **a}, è sbagliata,
come mostra la prossima sezione.
\end{affianco}

\subsection{\code{int **} non è una matrice}

Si potrebbe pensare di generalizzare quanto succede per gli array a una
dimensione, e dichiarare il parametro come puntatore a puntatori:

\begin{affianco}[0.5][c]
\begin{lstlisting}[style=c, name=sumpp.c]
#include <stdio.h>

int sum(int **a, int m, int n)
{
   int   i, j, sum = 0;

   for (i = 0; i < m; ++i)
      for (j = 0; j < n; ++j)
         sum += a[i][j];
   return sum;
}

int main()
{
   int A[3][5] = {{1, 2, 3, 4, 5}}, s;

   s = sum(A, 3, 5);
   printf("%d\n", s);
   return 0;
}
\end{lstlisting}
\risultato
\begin{console}
$ ./sumpp
Segmentation fault
\end{console}
\end{affianco}

\begin{console}
$ gcc -Wall -o sumpp sumpp.c
sumpp.c: In function (*@‘@*)main(*@’@*):
sumpp.c:17:12: warning: passing argument 1 of (*@‘@*)sum(*@’@*) from incompatible pointer type [-Wincompatible-pointer-types]
   17 |    s = sum(A, 3, 5);
      |            ^
      |            |
      |            int (*)[5]
sumpp.c:3:15: note: expected (*@‘@*)int **(*@’@*) but argument is of type (*@‘@*)int (*)[5](*@’@*)
\end{console}

Il compilatore avvisa, e il programma si blocca. Il motivo è che \code{A},
dichiarato \code{int A[3][5]}, è un puntatore ad \textbf{array} di 5 interi
(\code{int (*)[5]}), mentre un \code{int **} è un puntatore a
\textbf{puntatori}. Dentro \code{sum}, \code{a[i]} viene letto come un
indirizzo: ma nella memoria di \code{A} non ci sono indirizzi, ci sono gli
interi 1, 2, 3\dots, e usarli come indirizzi porta in zone di memoria
proibite (figura~\ref{fig:non-matrice}). Con un \code{int **}, invece,
\code{a[i]} è un puntatore che si può anche assegnare (\code{a[i] = \dots;}):
le due cose non sono equivalenti.

\begin{figure}[H]
  \centering
  \begin{tikzpicture}[font=\FiraCodeNL\scriptsize,
      cc/.style={mcella, minimum width=0.55cm}]
    \noindici
    \node[font=\Lato\scriptsize\bfseries] at (2.0,0.8) {\code{int A[3][5]}: 15 interi contigui};
    \celle[cc]{a}{(0,0)}{1,2,3,4,5,0,0,\ldots}
    \node[font=\Lato\scriptsize\bfseries] at (9.8,0.8) {\code{int **a}: puntatori a righe};
    \memoria[minimum width=0.9cm]{p}{(7.6,0)}{a[0]/{}/, a[1]/{}/, a[2]/{}/}
    \foreach \r in {1,2,3} {
      \celle[cc]{r\r}{(9.6,-\r*0.52+0.52)}{\mvuota,\mvuota,\mvuota,\mvuota,\mvuota}
      \riferimento{p-\r}{-- (r\r-0.west)}
    }
    \node[mnota, anchor=north, align=center, text=merr] at (2.2,-0.45)
      {se \code{sum} legge queste celle come indirizzi,\\“salta” all'indirizzo 1, 2\dots: memoria proibita};
  \end{tikzpicture}
  \caption{Un array a due dimensioni non contiene puntatori; un
           \code{int **} sì.}
  \label{fig:non-matrice}
\end{figure}

\subsection{Passare l'indirizzo base}

Un modo per scrivere una funzione valida per qualunque dimensione è passarle
l'\textbf{indirizzo base} e le dimensioni, e applicare a mano la storage
mapping function:

\begin{affianco}[0.5][c]
\begin{lstlisting}[style=c]
int sum_base(int *a, int m, int n)
{
   int   i, j, sum = 0;

   for (i = 0; i < m; ++i)
      for (j = 0; j < n; ++j)
         sum += *(a + n*i + j);
   return sum;
}
\end{lstlisting}
\risultato
Si chiama con \code{sum\_base(\&A[0][0], 3, 5)}. Con una matrice
\code{A[3][5]} tutta di 1, sia \code{sum(A, 3)} sia
\code{sum\_base(\&A[0][0], 3, 5)} restituiscono 15 (verificato). Funziona,
ma si perde la comoda scrittura \code{a[i][j]}.
\end{affianco}

\section{Le matrici dinamiche}

\subsection{Un array di righe}

Gli array a più dimensioni, insomma, non sono molto flessibili: non si
possono passare alle funzioni senza perdere in generalità (dimensioni fisse)
o in semplicità (indirizzo base). Sarebbe più comodo poter scrivere funzioni
come \code{int f(int **a, int m, int n);}, a cui si passa la matrice e le
sue dimensioni. Si può, allocando la matrice dinamicamente come un
\textbf{array di puntatori alle righe}:

\begin{affianco}[0.5][c]
\begin{lstlisting}[style=c, name=matrici.c]
double **matrix_alloc(int m, int n)
{
   int i;
   double **p;

   p = malloc(m * sizeof(double *));
   for (i = 0; i < m; i++)
      p[i] = malloc(n * sizeof(double));
   return p;
}

void free_matrix(double **a, int m)
{
   int i;

   for (i = 0; i < m; i++)
      free(a[i]);
   free(a);
}
\end{lstlisting}
\risultato
\begin{itemize}
  \item Prima si alloca l'array di \code{m} puntatori (ognuno occupa
        \code{sizeof(double *)}), poi ogni riga, di \code{n} \code{double}.
  \item All'elemento si accede con la comoda scrittura \code{a[i][j]}:
        \code{a[i]} è il puntatore alla riga $i$, e \code{[j]} sceglie
        l'elemento nella riga.
  \item La memoria della matrice \textbf{non è contigua}: ogni riga sta
        dove \code{malloc} l'ha messa.
  \item Si libera con una \code{free} per ogni riga, e infine quella
        dell'array dei puntatori.
\end{itemize}
\end{affianco}

\subsection{Memoria contigua}

Una variante alloca tutti gli elementi in un'unica zona contigua, e fa
puntare le righe nei punti giusti: \code{a[0]} all'inizio, e ogni riga $n$
elementi dopo la precedente (figura~\ref{fig:matrici-dinamiche}).

\begin{affianco}[0.5][c]
\begin{lstlisting}[style=c]
double **matrix_alloc2(int m, int n)
{
   int i;
   double **a;

   a = malloc(m * sizeof(double *));
   a[0] = malloc(m * n * sizeof(double));
   for (i = 1; i < m; i++)
      a[i] = a[i-1] + n;
   return a;
}

void free_matrix2(double **a)
{
   free(a[0]);
   free(a);
}
\end{lstlisting}
\risultato
\begin{console}
$ ./matrici
a[2][3] = 23   b[2][3] = 23
b: riga 1 inizia 4 double dopo la riga 0
\end{console}
\wtnota{\code{a} è allocata riga per riga, \code{b} in modo contiguo; entrambe
sono $3 \times 4$, con \code{[i][j]} uguale a $10 i + j$. Compilato con
\code{-fsanitize=address}, il programma non dà errori.}
\end{affianco}

Nelle slide la seconda funzione di liberazione è dichiarata
\code{double free\_matrix(double **a)}: non restituisce nulla, quindi il
tipo giusto è \code{void} (con \code{double}, \code{gcc -Wall} avvisa
\textit{control reaches end of non-void function}).

\begin{figure}[H]
  \centering
  \begin{tikzpicture}[font=\FiraCodeNL\scriptsize,
      cc/.style={mcella, minimum width=0.55cm}]
    \noindici
    \node[font=\Lato\scriptsize\bfseries] at (2.3,0.75) {riga per riga};
    \memoria[minimum width=0.9cm]{p}{(0,0)}{a[0]/{}/, a[1]/{}/, a[2]/{}/}
    \foreach \r/\y in {1/0.25, 2/-0.55, 3/-1.35} {
      \celle[cc]{q\r}{(2.0,\y)}{{},{},{},{}}
      \riferimento{p-\r}{-- (q\r-0.west)}
    }
    \node[mnota, anchor=north] at (2.6,-1.75) {tre zone separate};
    \begin{scope}[xshift=7.6cm]
      \node[font=\Lato\scriptsize\bfseries] at (2.6,0.75) {memoria contigua};
      \memoria[minimum width=0.9cm]{s}{(0,0)}{b[0]/{}/, b[1]/{}/, b[2]/{}/}
      \celle[cc]{z}{(2.0,-1.3)}{{},{},{},{},{},{},{},{},{},{},{},{}}
      \riferimento{s-1}{to[out=0, in=90] (z-0.north)}
      \riferimento{s-2}{to[out=0, in=90] (z-4.north)}
      \riferimento{s-3}{to[out=0, in=90] (z-8.north)}
      \node[mnota, anchor=north] at ([yshift=-0.15cm]z-6.south) {un'unica zona: \code{b[i] = b[i-1] + n}};
    \end{scope}
  \end{tikzpicture}
  \caption{Due modi di allocare una matrice $3 \times 4$.}
  \label{fig:matrici-dinamiche}
\end{figure}

\section{Gli array a lunghezza variabile}

\subsection{Le dimensioni decise durante l'esecuzione}

Il C99 ha introdotto gli \term{array a lunghezza variabile} (\textit{variable
length array}, VLA): array automatici la cui dimensione è un'espressione
calcolata durante l'esecuzione, per esempio un valore letto da tastiera.
Ecco una versione ridotta dell'esempio delle slide (che a sua volta viene
dal libro di Deitel):

\begin{affianco}[0.55][c]
\begin{lstlisting}[style=c, name=vla.c]
#include <stdio.h>

void print2DArray(int row, int col, int array[row][col]);

int main(void)
{
   int row, col;

   printf("Righe e colonne: ");
   scanf("%d %d", &row, &col);

   int m[row][col];                /* un VLA */

   for (int i = 0; i < row; ++i)
      for (int j = 0; j < col; ++j)
         m[i][j] = i + j;
   printf("sizeof(m) = %zu byte\n", sizeof(m));
   print2DArray(row, col, m);
   return 0;
}

void print2DArray(int row, int col, int array[row][col])
{
   for (int i = 0; i < row; ++i) {
      for (int j = 0; j < col; ++j)
         printf("%5d", array[i][j]);
      puts("");
   }
}
\end{lstlisting}
\risultato
\begin{console}
$ ./vla
Righe e colonne: <<2 5>>
sizeof(m) = 40 byte
    0    1    2    3    4
    1    2    3    4    5
$ ./vla
Righe e colonne: <<4 3>>
sizeof(m) = 48 byte
    0    1    2
    1    2    3
    2    3    4
    3    4    5
\end{console}
\end{affianco}

\begin{itemize}
  \item La dichiarazione \code{int m[row][col];} viene dopo la lettura delle
        dimensioni: in C99 le dichiarazioni si possono mettere in mezzo alle
        istruzioni, e così le variabili dei \code{for}
        (\code{for (int i = 0; \dots)}).
  \item \code{sizeof(m)} viene calcolato durante l'esecuzione: $2 \cdot 5
        \cdot 4 = 40$ byte la prima volta, $4 \cdot 3 \cdot 4 = 48$ la
        seconda.
  \item Nel parametro \code{int array[row][col]} le dimensioni sono altri
        parametri, dichiarati prima: la funzione va bene per matrici di
        qualunque dimensione.
\end{itemize}

\subsection{Pro e contro}

I VLA sono comodi, ma hanno un difetto: sono allocati nello
\textbf{stack} (capitolo~\ref{chap:puntatori}), che è piccolo rispetto allo
heap, e la loro dimensione non è nota in compilazione. Con dimensioni grandi
c'è il rischio concreto di esaurire lo stack: con \code{2000 2000} la matrice
occupa 16 MB, il doppio degli 8 MB che Linux assegna di solito allo stack, e
il programma viene fermato con \textit{Segmentation fault}. I VLA sono
quindi utilizzabili per piccole dimensioni; per le grandi si usa
\code{malloc}. È anche il parere di Linus Torvalds, il creatore di Linux,
riportato nel file \code{input.txt} dell'ordinamento delle parole: i VLA
\textit{bring nothing of positive value}, meglio un array fisso
sufficientemente grande o l'allocazione nello heap.

\insintesi{
  Una \textbf{stringa} è un array di caratteri terminato da
  \code{'\textbackslash 0'}, che va contato nella dimensione. Una stringa
  costante è un puntatore a memoria di sola lettura; \code{\%s} stampa fino
  al terminatore. \code{const} impedisce di modificare una variabile o il
  dato puntato, e il compilatore lo verifica. \code{ctype.h} classifica e
  trasforma i caratteri, \code{string.h} copia, concatena, confronta e cerca
  nelle stringhe, senza controllare lo spazio: per una copia servono
  \code{strlen + 1} byte. Un \textbf{array di puntatori} raccoglie stringhe
  di lunghezza diversa, e si ordina scambiando i puntatori. Gli array a
  \textbf{più dimensioni} stanno in memoria riga dopo riga (storage mapping
  function \code{\&a[0][0] + N*i + j}); come parametri vanno dichiarate
  tutte le dimensioni tranne la prima, e un \code{int **} non è una matrice.
  Le matrici dinamiche sono array di puntatori alle righe; i VLA hanno
  dimensioni decise in esecuzione, ma stanno nello stack.}

\section{Riferimenti}

\begin{itemize}
  \item \textbf{A. Kelley, I. Pohl}, \textit{A Book on C}: capitolo~6
        (stringhe, array di puntatori, array a più dimensioni).
  \item \textbf{H. M. Deitel, P. J. Deitel}, \textit{C: How to Program}:
        array a lunghezza variabile.
  \item \textbf{Null-terminated byte strings} --- cppreference.com. \\
        \urlx{https://en.cppreference.com/w/c/string/byte}
\end{itemize}
